% !TeX root = ../教材题库_全部条目_转写工作稿.tex
% Source: XJ-B1-C06-PR-P247-Y0217
% 第6章 / 6.4 用样本估计总体 / 教材 p.247 / PDF 第254页
% OCR未校对：公式与文字均可能有误；未完成数学TeX转写；禁止直接发给学生
% 原题及相邻原解截图（核对用）：../题目原图/XJ-B1-C06-PR-P247-Y0217-01.webp
% 本文件只在转写完成后作为题干使用；OCR 条目尚不是可靠题干。
\begin{bankitem}{p.247 练习 第1题}{第6章 / 6.4 用样本估计总体 / 教材 p.247 / PDF 第254页}{XJ-B1-C06-PR-P247-Y0217}{OCR未校对：公式与文字均可能有误；未完成数学TeX转写；禁止直接发给学生}
\begin{ocrdraft}
[教材p.247 / PDF第254页]
\par
1.一家公司准备盖大楼，有两块土地可供征用，但两块土地都崎岖不平，
\par
需要平整．现对每块土地确定房基基准高度，然后在两块土地上分别适当地另取
\par
10点，用水平仪测得各点相对基准的标高（单位：cm）如下表所示：
\par
2
\par
3
\par
4
\par
5
\par
6
\par
7
\par
8
\par
9
\par
47
\par
一26
\par
135
\par
84
\par
一61
\par
一38
\par
76
\par
92
\par
乙
\par
74
\par
120
\par
100
\par
一70
\par
一44
\par
95
\par
63
\par
一50
\par
57
\par
一25
\par
问：哪一块土地较容易平整？
\end{ocrdraft}
\end{bankitem}
